\section{An unconditional Mordell--Weil family in iPRA}
\label{sec:unconditional}

We now discharge all arithmetic hypotheses for a concrete elliptic curve,
and identify its group rather than merely bounding its rank.
The proof does not use the supplied-cover theorem, a search cutoff,
Selmer groups, or a theorem about all elliptic curves.
Its engine is a completely bounded version of Fermat's square-area descent.
Related arguments were present in the inspected ProveIt Diophantine and
Faltings packages \cite{proveit-geometry,proveit-fermat};
the complete argument needed here is included.

\subsection{Primitive square-area triangles}

\begin{definition}
A square-area triangle is a positive integer tuple $(a,b,c,s)$ satisfying
\[
       a^2+b^2=c^2,\qquad ab=2s^2.
 \tag{6.1}\label{eq:triangle}
\]
It is primitive if $\gcd(a,b)=1$.
\end{definition}

\begin{lemma}[Normalization]\label{lem:normal-triangle}
Every square-area triangle yields a primitive square-area triangle
$(A,B,C,S_0)$ with $C\le c$, by a specified primitive recursive operation.
\end{lemma}
\begin{proof}
Let $g=\gcd(a,b)>0$ and put $A=a/g,B=b/g$.
Then $g^2\mid c^2$, so Lemma~\ref{lem:squares} gives $g\mid c$.
Set $C=c/g$. The pair $A,B$ is coprime and
$A^2+B^2=C^2$.

Exactly one of $A,B$ is even. They cannot both be even by coprimality,
and cannot both be odd because the only square residues modulo $4$ are
$0,1$, whereas two odd squares sum to $2$. Therefore $K=AB/2$ is an
integer. From $ab=2s^2$ we get $s^2=g^2K$.
Again Lemma~\ref{lem:squares} gives $g\mid s$. Put $S_0=s/g$,
cancel $g^2$, and obtain $AB=2S_0^2$.
All four entries are positive and $C=c/g\le c$.
\end{proof}

\begin{lemma}[Bounded Pythagorean parametrization]\label{lem:pyth}
For a primitive positive triangle $A^2+B^2=C^2$, exchange the legs
if necessary so that $A$ is even. There are computable integers $m>n>0$
with
\[
 A=2mn,\qquad B=m^2-n^2,\qquad C=m^2+n^2,
 \tag{6.2}\label{eq:pyth}
\]
where $\gcd(m,n)=1$ and $m,n$ have opposite parity.
The required square-root scans have bound $C$.
\end{lemma}
\begin{proof}
The parity argument makes $B,C$ odd. The pair $B,C$ is coprime:
a common divisor $d$ would satisfy $d^2\mid C^2-B^2=A^2$,
whence $d\mid A$, so $d=1$. Since $C>B>0$, the numbers
\[
              U=(C+B)/2,\qquad V=(C-B)/2
\]
are positive, $U>V$, and $UV=(A/2)^2$.
A common divisor of $U,V$ divides $C$ and $B$, so they are coprime.
Square splitting gives $U=m^2,V=n^2,A/2=mn$.
Their sum and difference give \eqref{eq:pyth}.

A common divisor of $m,n$ would divide $U,V$; hence they are coprime.
Their squares sum to the odd number $C$, so they have opposite parity.
The nonnegative roots are unique, since
$t^2-r^2=(t-r)(t+r)>0$ when $0\le r<t$.
Both radicands are at most $C$, and their roots can therefore be found
by scans of $0,\ldots,C$. Every division by two has an even numerator.
\end{proof}

\subsection{The explicit descending transformation}

\begin{lemma}[A strictly smaller square-area triangle]\label{lem:fermat-step}
A primitive square-area triangle $(A,B,C,S_0)$ gives a primitive
square-area triangle $(A',B',C',S'_0)$ with $0<C'<C$.
The transformation is a fixed total primitive recursive function after
assigning a default output to invalid inputs.
\end{lemma}
\begin{proof}
Use the preceding parametrization. The area equation becomes
\[
             mn(m+n)(m-n)=S_0^2.
 \tag{6.3}\label{eq:fourfactors}
\]
The four factors are positive and pairwise coprime.
For pairs involving $m$ or $n$, subtract multiples to reduce a common
divisor to one of $m,n$. A common divisor of $m+n,m-n$ divides
$2m,2n$, and hence divides $2$ by a Bezout equation for $m,n$.
Both $m+n,m-n$ are odd, so this divisor is $1$.

Repeated square splitting now supplies positive integers $r,s,t,u$ with
\[
        m=r^2,\quad n=s^2,\quad m+n=t^2,\quad m-n=u^2.
 \tag{6.4}\label{eq:fourroots}
\]
In particular $t>u>0$, and $t,u$ are odd and coprime.
The integer $s$ is even. To prove this, if $m$ were even and $n$ odd,
their square forms would give $m\equiv0,n\equiv1\pmod4$,
so $m-n=u^2\equiv3\pmod4$, a contradiction.
Thus $n$ is even, and its square root $s$ is even.

Define
\[
        A'=(t+u)/2,\quad B'=(t-u)/2,\quad C'=r,\quad S'_0=s/2.
 \tag{6.5}\label{eq:fermat-map}
\]
These are positive integers. Expanding gives
\[
 (A')^2+(B')^2=(t^2+u^2)/2=m=r^2,
 \qquad
 A'B'=(t^2-u^2)/4=n/2=2(s/2)^2.
 \tag{6.6}
\]
A common divisor of $A',B'$ divides their sum $t$ and difference $u$,
so the new triangle is primitive. Finally
\[
                 0<C'=r\le r^2=m<m^2+n^2=C.
 \tag{6.7}
\]

To specify a total program, first test the primitive triangle conditions.
On failure return $(0,0,0,0)$ with a failure bit. Otherwise order the legs,
find $m,n$ as in Lemma~\ref{lem:pyth}, find the four roots in
\eqref{eq:fourroots}, and return \eqref{eq:fermat-map}.
Every root scan may use the original $C$:
$m,n,m-n,m+n\le m^2+n^2=C$.
The proof establishes that every demanded root exists and every exact
division succeeds in the valid branch. An invalid root search also has a
specified default output.
\end{proof}

\begin{theorem}[Square-area impossibility in iPRA]\label{thm:no-square-area}
There are no positive integers satisfying \eqref{eq:triangle}.
\end{theorem}
\begin{proof}
Suppose a solution is given and normalize it using
Lemma~\ref{lem:normal-triangle}. Let $q$ code the resulting primitive
triangle, $\Good(q)$ be its primitive recursive Boolean validity
predicate, and $\rho(q)$ its hypotenuse. Let $D$ be the total map just
constructed, and define $q_0=q,q_{j+1}=D(q_j)$ by primitive recursion.
Under the assumption $\Good(q)=1$, quantifier-free induction proves
\[
           \Good(q_j)=1\quad\wedge\quad \rho(q_j)+j\le\rho(q).
 \tag{6.8}\label{eq:fermat-invariant}
\]
The base is immediate. The step follows from
$\Good(D(q_j))=1$ and $\rho(D(q_j))+1\le\rho(q_j)$.
All relations in the invariant are Boolean tests or equations of the
fixed coded functions.

Instantiate this invariant at the explicit term $j=\rho(q)$.
It implies $\rho(q_j)=0$, while validity requires $\rho(q_j)>0$.
This contradiction discharges the assumed solution.
The proof runs for exactly the initial hypotenuse; it chooses neither a
least counterexample nor an infinite descending sequence.
Proving the negation of a solution by deriving falsity from that solution
is intuitionistically valid.
\end{proof}

\subsection{All rational points of \texorpdfstring{$y^2=x^3-x$}{y2=x3-x}}

\begin{theorem}[Exact rational point classification]\label{thm:e1}
In the rational codes defined above, iPRA proves
\[
 E_1(\Q)=\{O,(0,0),(1,0),(-1,0)\},
 \qquad E_1:\ y^2=x^3-x.
 \tag{6.9}\label{eq:e1points}
\]
In particular every rational point has $x$-height $1$.
\end{theorem}
\begin{proof}
The projective point at infinity is $O$, so consider an affine point
$(x,y)$. Rational equality is decidable. In the branch $y\ne0$,
the equation excludes $x=0,1,-1$. Define the positive rational numbers
\[
 \alpha=\left|\frac{x^2-1}{y}\right|,
 \qquad
 \beta=\left|\frac{2x}{y}\right|,
 \qquad
 \gamma=\left|\frac{x^2+1}{y}\right|.
 \tag{6.10}\label{eq:triangle-map}
\]
They are strictly positive by the exclusions and $x^2+1>0$.
The polynomial identity
$(x^2-1)^2+4x^2=(x^2+1)^2$ gives
$\alpha^2+\beta^2=\gamma^2$. The curve equation gives
\[
 \frac{\alpha\beta}{2}
   =\left|\frac{x(x^2-1)}{y^2}\right|=1.
 \tag{6.11}
\]
Let $d>0$ be the product of the positive denominators of
$\alpha,\beta,\gamma$. Then
$a=d\alpha,b=d\beta,c=d\gamma$ are positive integers and
$a^2+b^2=c^2$, $ab=2d^2$. This contradicts
Theorem~\ref{thm:no-square-area}, with $s=d$.

Thus the decidable case split on $y=0$ leaves only $y=0$.
Then $x(x-1)(x+1)=0$. The integral-domain property of the rational
arithmetic, or finite case distinction on these three factors,
gives $x=0,1,-1$. Direct substitution verifies all three points.
The list is pairwise distinct, and their heights, including $O$, are $1$.

Every operation in this reduction is normalization, exact rational
arithmetic, sign comparison, or the bounded triangle procedure.
The nonexistence proof was already given with its explicit induction
formula. No empirical point enumeration is used to conclude completeness.
\end{proof}

\subsection{The group, generators, relations and coordinate function}

Write $T=(0,0)$, $U=(1,0)$, $V=(-1,0)$. Each has negative equal to
itself, so its double is $O$. The line through any two distinct affine
points in this list is $y=0$. Formula \eqref{eq:add} shows that the sum
is the remaining affine point. Thus the complete table is
\[
\begin{array}{c|cccc}
 +&O&T&U&V\\ \hline
 O&O&T&U&V\\
 T&T&O&V&U\\
 U&U&V&O&T\\
 V&V&U&T&O
\end{array}
\tag{6.12}
\]
Identify these labels respectively with
$(0,0),(1,0),(0,1),(1,1)\in\{0,1\}^2$.
The table is componentwise addition modulo two. This proves all group
laws, the isomorphism
\[
                     E_1(\Q)\cong(\Z/2\Z)^2,
 \tag{6.13}
\]
and the presentation with generators $T,U$ and relations
$2T=2U=0$, $T+U=U+T$. There is no additional relation between the two
bits: their four sums are the four distinct points.

The coordinate algorithm tests equality with the four listed points
and returns the associated pair of bits. By Theorem~\ref{thm:e1} the
successful branch always exists for a valid point; finite search supplies
the witnesses. It therefore proves a genuine constructive finite-generation
statement, including exact coordinates and relations. The finite table
is separately checked in \code{lean/E1FiniteTable.lean}; the arithmetic
surjectivity onto that table is the preceding written proof.

\subsection{Uniform rational scalings and finite products}

\begin{theorem}[An unconditional family in every finite dimension]
\label{thm:products}
For every rational $u\ne0$ let $E_u:y^2=x^3-u^4x$. Then iPRA proves
\[
 E_u(\Q)=\{O,(0,0),(u^2,0),(-u^2,0)\},
 \qquad E_u(\Q)\cong(\Z/2\Z)^2.
 \tag{6.14}
\]
For every finite input list $u_1,\ldots,u_r\in\Q^\times$, the rational
group of the product has the explicit presentation
\[
 \left(\prod_{j=1}^rE_{u_j}\right)(\Q)
       \cong(\Z/2\Z)^{2r}.
 \tag{6.15}
\]
It has exactly $4^r$ points, rank zero, and a primitive recursive
$2r$-bit coordinate map. The $2r$ coordinate-axis points are generators.
For $r=0$ this means the trivial group and an empty generator list.
\end{theorem}
\begin{proof}
The mutually inverse rational maps
\[
 (x,y)\longmapsto(u^2x,u^3y),\qquad
 (X,Y)\longmapsto(X/u^2,Y/u^3)
 \tag{6.16}
\]
fix $O$ and carry the respective curve equations to one another.
Their denominators are nonzero because $u\ne0$.
Theorem~\ref{thm:e1} therefore proves the displayed list for $E_u$.
Its three nonidentity ordinates are zero; applying
\eqref{eq:add} gives the same four-element table. Thus the
coordinate bijection is a group isomorphism, without needing any
separate abstract theorem about coordinate changes.

A rational point of the product is a tuple of points in its factors.
Apply the four-case test independently to every coordinate, obtaining
$2r$ bits. Conversely any $2r$ bits give exactly one such tuple.
Both transformations are bounded list folds, inverse by their
coordinatewise identities. Addition is bitwise modulo two, so the
coordinate-axis tuples generate and the number of points is $2^{2r}=4^r$.
Induction on the finite input list can be expressed by a Boolean test
of all already processed coordinates together with the accumulated
integer count; it is quantifier-free after list coding.

The ordinary geometric description of this product as an abelian variety
of dimension $r$ is background terminology. The iPRA assertion proved
here is the explicit statement about rational tuple codes, their
membership tests, group table and coordinate map; it does not quantify
over schemes.
\end{proof}

The models in this family are rationally isomorphic to powers of the same
elliptic curve. This is an unconditional infinite family of input
presentations, not a claim that arbitrary twists, positive-rank curves,
or arbitrary number fields have been settled.
